\section{Bounded Power-Graph Chromatic Number Under Bounded Degree}

\begin{theorem}[Bounded $\chi_R$ for Bounded Degree Graphs]
Let $G_n$ be a family of graphs with maximum degree $\Delta$ independent of $n$.
Fix $R \ge 1$ independent of $n$.
Then the $2R$-power graph $G_n^{(2R)}$ satisfies

\[
\chi(G_n^{(2R)}) \le 1 + \Delta \sum_{i=0}^{2R-1} (\Delta-1)^i
\]

In particular, $\chi(G_n^{(2R)}) = O_{\Delta,R}(1)$.
\end{theorem}

\begin{proof}[Sketch]
For any vertex $v$, the number of vertices within distance $2R$ is bounded by

\[
1 + \Delta \sum_{i=0}^{2R-1} (\Delta-1)^i.
\]

Thus $G_n^{(2R)}$ has maximum degree bounded by a constant depending only on $\Delta$ and $R$.
Greedy coloring yields $\chi(G_n^{(2R)}) \le \Delta(G_n^{(2R)}) + 1$.
\end{proof}

\section{Entropy-Depth Consequence}

\begin{corollary}[Constant Depth Under Local Bounded Degree]
Under the assumptions above, any refinement process limited to $R$-local activation satisfies

\[
T \ge \chi(G_n^{(2R)}) = O_{\Delta,R}(1).
\]

Thus bounded locality and bounded degree alone do not force superconstant entropy depth.
\end{corollary}

\paragraph{Implication.}
Linear entropy depth requires violation of at least one:
\begin{itemize}
\item Unbounded degree
\item $R$ growing with $n$
\item Diameter collapse
\item Nonlocal transcript accumulation
\end{itemize}
